% GPT-6 Astra Pro, /root/replacement, 2026-11-22.
% Actual author proof; effort setting not exposed.
\section{HD-F1-R05: abelian-support sums fail the whole-character gate}
\label{sec:f1-abelian-support}

Let $(A,K,\theta)$ be a polarized complex abelian variety of
Weil type of dimension $2n$ and signature $(n,n)$, with $n\ge2$.
Let
\[
 j_\nu:B_\nu+a_\nu\hookrightarrow A
\]
be inclusions of translates of $n$-dimensional abelian
subvarieties. Let $E_\nu$ be coherent sheaves of positive
generic rank on these supports, and let $b_\nu\in\mathbf Z$.
Consider the actual nonempty direct sum
\[
 F=\bigoplus_{\nu=1}^s j_{\nu,*}E_\nu[b_\nu].
 \tag{F1R5.1}
\]
For every such object satisfying the whole-character identity
\[
 e^{B_0}\operatorname{ch}(F)
       =P(\theta)+\lambda t,
 \qquad B_0\in H^{1,1}(A,\mathbf Q),\quad
 t\in W_K(A),\quad\lambda t\ne0,
 \tag{F1R5.2}
\]
the full categorical trace has a nonzero kernel on
$\operatorname{Ext}^2(F,F)^G$ for every actual coherent action
of a faithful finite point subgroup $G\subset A\times\widehat A$.
Thus no such object is weakly $G$-semiregular.

The polynomial $P$ is arbitrary. The hypotheses allow repeated
supports, all intersections, arbitrary shifts, and arbitrary
mixing in the specified coherent linearization. In particular
they include sums of transverse graph structure sheaves with
character $c\theta^n+\lambda t$. No assertion is made for a
complex with new differentials between the displayed summands.
This extends the actual-support result HD-N21 from its pure
Weil character to the full S-14 character.

\paragraph{1. Every proper middle support moves to first order.}
Let $V=H_1(A,\mathbf R)$ and use the action of $K$ to regard
it as a complex vector space of dimension $2n$. Write $J$ for
this scalar complex structure and $I$ for the complex structure
defining $A$. The Hermitian form associated to $\theta$ has
signature $(n,n)$. Choosing orthonormal bases in its positive
and negative Hodge subspaces gives
\[
 V=\mathbf C^n\oplus\mathbf C^n,\quad
 h=\operatorname{diag}(1_n,-1_n),\quad
 I=\operatorname{diag}(i1_n,-i1_n).
 \tag{F1R5.3}
\]
For every complex matrix $Z\in\operatorname{Mat}_n(\mathbf C)$,
put
\[
 P_Z=\begin{pmatrix}0&Z\\Z^*&0\end{pmatrix}.
 \tag{F1R5.4}
\]
It satisfies $P_Z^*h+hP_Z=0$ and anticommutes with $I$.
Therefore $I_s=e^{sP_Z}Ie^{-sP_Z}$ is an actual local family
of polarized $K$-linear complex structures of signature $(n,n)$.
The lattice and polarization are fixed. Positivity is preserved
because the conjugating operator is $h$-unitary. Its derivative
is
\[
 [P_Z,I]=P_{-2iZ}.
 \tag{F1R5.5}
\]
Consequently the real tangent directions in this family run
over all the matrices $P_Z$.

Here is the complete algebra-generation calculation. For
$i\ne k$, fixed $j$, and $\zeta\in\mathbf C$, the elementary
matrices satisfy
\[
 P_{\zeta E_{ij}}P_{E_{kj}}
          =\operatorname{diag}(\zeta E_{ik},0).
 \tag{F1R5.6}
\]
Products with the reverse matrix unit give
$\operatorname{diag}(\zeta E_{ii},0)$. Since $n\ge2$, these
identities produce every matrix in the upper diagonal block.
For fixed $i$ and $j\ne l$, the analogous identity is
\[
 P_{E_{ij}}P_{\zeta E_{il}}
          =\operatorname{diag}(0,\zeta E_{jl}).
 \tag{F1R5.7}
\]
It produces the entire lower diagonal block. Multiplying
$P_Z$ by the two resulting block projectors then produces
every upper and lower off-diagonal matrix. The real associative
algebra generated by all $P_Z$ is thus
$\operatorname{Mat}_{2n}(\mathbf C)$, including scalar $J$.

Fix one abelian support $B$. Its real tangent space
$W=H_1(B,\mathbf R)\subset V$ is a nonzero proper $I$-stable
subspace. It cannot be stable under every $P_Z$: stability
would make it a module for the full displayed matrix algebra,
and first for $J$, so a complex subspace stable under all
complex endomorphisms of $V$. Such a subspace is zero or $V$.
Thus some polarized Weil-family tangent fails to preserve $W$.

To identify exactly the corresponding cohomology map, put
$M=T_0^{1,0}A$, $L=T_0^{1,0}B$, and let
$\kappa\in\operatorname{Hom}(\overline M,M)
 =H^1(A,T_A)$ be the Kodaira--Spencer class of that tangent.
The infinitesimal action on forms is the derivation
$D_\kappa$ of Hodge degree $(-1,1)$.
Choose bases $\alpha_1,\ldots,\alpha_n$ of
$\operatorname{Ann}(L)\subset M^*$ and their conjugates.
The quotient map $q:A\to A/B$ gives
\[
 [B]=q^*[0]
      =c\,\alpha_1\wedge\cdots\wedge\alpha_n
                  \wedge\overline\alpha_1\wedge\cdots
                          \wedge\overline\alpha_n,
 \qquad c\ne0.
 \tag{F1R5.8}
\]
The independent $(n-1)$-fold holomorphic factors in its
derivative show that
\[
 D_\kappa[B]=0
 \quad\Longleftrightarrow\quad
 D_\kappa\alpha_i\in\operatorname{Ann}(\overline L)
          \text{ for every }i
 \quad\Longleftrightarrow\quad
 \kappa(\overline L)\subset L.
 \tag{F1R5.9}
\]
Indeed wedge with the antiholomorphic normal volume has kernel
exactly its $n$-dimensional span. The last condition and its
conjugate say precisely that the real derivative of $I$
preserves $W$. Our chosen tangent therefore satisfies
\[
                         D_\kappa[B]\ne0.
 \tag{F1R5.10}
\]
This establishes the required direction by actual period,
restriction and Gysin maps, without a generic-simplicity claim.

\paragraph{2. The full character gives a trace-zero operator.}
Every tangent above preserves $\theta$, and the two Weil
determinant lines remain of type $(n,n)$ throughout the family.
In the $K$-eigenform decomposition each determinant already
contains all $n$ antiholomorphic variables of its eigenblock;
$D_\kappa$ preserves that block and its replacement of a
holomorphic variable therefore wedges a repeated
antiholomorphic variable. Hence, directly,
\[
 D_\kappa\theta=0,\qquad D_\kappa t=0,
 \qquad D_\kappa(P(\theta)+\lambda t)=0.
 \tag{F1R5.11}
\]

Set $\eta=D_\kappa B_0\in H^2(A,\mathcal O_A)$ and let
\[
 a=\kappa+\eta\in
 H^1(A,T_A)\oplus H^2(A,\mathcal O_A)\subset HH^2(A).
 \tag{F1R5.12}
\]
Its action on forms is $D_\kappa+\eta\wedge(-)$.
The exact exponential identity is
\[
 (D_\kappa+\eta\wedge)
          \bigl(e^{-B_0}w\bigr)=e^{-B_0}D_\kappa w.
 \tag{F1R5.13}
\]
It follows by differentiating the exponential; no line bundle
with Chern class $B_0$ is needed or asserted. Applying it to
$w=P(\theta)+\lambda t$ proves
\[
 a\,\operatorname{ch}(F)=0.
 \tag{F1R5.14}
\]
Under the actual Hochschild characteristic morphism and full
Atiyah trace identity, this says
\[
 \sigma_F^2\bigl(\operatorname{ev}_F(a)\bigr)=0.
 \tag{F1R5.15}
\]
The trace/evaluation identity is the same published one used
in HD-N21: Buchweitz--Flenner, Adv. Math.217 (2008),
Proposition6.2.1, Corollary6.3.2 and Theorem6.4.2,
author version arXiv:math/0606730, pp.29--31. Its HKR
Todd correction is one on an abelian variety.

\paragraph{3. The operator is nonzero on an actual summand.}
For the support chosen in paragraph1, write its positive
generic rank as $r_B$. GRR for its closed immersion, whose
source and target tangent bundles are trivial, gives
\[
 \operatorname{ch}_k(j_*E)=0\ (k<n),\qquad
 \operatorname{ch}_n(j_*E)=r_B[B].
 \tag{F1R5.16}
\]
The component of $a\operatorname{ch}(j_*E)$ in
$H^{n+1}(A,\Omega_A^{n-1})$ is exactly
\[
                            r_BD_\kappa[B]\ne0.
 \tag{F1R5.17}
\]
The $\eta$-term has holomorphic degree at least $n$;
the $D_\kappa$-terms on higher character components do too.
Thus none can cancel (F1R5.17). The trace/evaluation identity
on this summand proves $\operatorname{ev}_{j_*E}(a)\ne0$.
Shifting only introduces a nonzero sign in its trace.
Naturality of the characteristic morphism makes its value
on an actual direct sum diagonal. The nonzero summand just
proved therefore gives
\[
                     \operatorname{ev}_F(a)\ne0.
 \tag{F1R5.18}
\]
This step does not assume that the cross-Ext groups vanish.

\paragraph{4. Every permitted action retains the kernel, with an exact weak comparison.}
A point functor $t_x^*(-)\otimes P$, $P\in\operatorname{Pic}^0(A)$,
acts identically on $HH^*(A)$: invariant forms and vector
fields are translation invariant, and the tensor shear is
zero because the Atiyah class of $P$ vanishes. Naturality
of evaluation consequently makes (F1R5.18) invariant under
every specified coherent linearization, including those which
mix summands.

For completeness the necessary weak comparison is an
isomorphism, rather than the converse of a sufficient test.
Let $\mathcal C=\operatorname{Perf}(A)$. Perry's published
coefficient-sector decomposition, Math. Z.297 (2021),
Theorem4.4, gives
\[
 HH^*(\mathcal C^G)
   =\left(\bigoplus_{g\in G}HH^*(\mathcal C,\Phi_g)\right)^G.
 \tag{F1R5.19}
\]
The graph of a nonzero translation is disjoint from the
diagonal, so that sector vanishes. A nontrivial pure tensor
sector is $\bigoplus H^q(A,P)\otimes\bigwedge^pT_0A=0$.
For torsion $P$ this vanishing follows by the connected
isogeny trivializing $P$: deck translations act trivially
on all $H^q(\mathcal O)$, leaving only the trivial character.
Faithfulness leaves the identity as the sole surviving sector.
All point functors preserve the invariant volume, so
the invariant category is Calabi--Yau of dimension $2n$.
Thus its Hochschild homology and that of $A$ have equal
dimensions in every degree. Actual induction and forgetting
satisfy
\[
 \mathrm{Forg}_*\mathrm{Ind}_*=|G|\operatorname{id}_{HH_*(A)}.
 \tag{F1R5.20}
\]
Their displayed surjective forgetful map is therefore an
isomorphism. Moreover
$\operatorname{Ext}^2_{\mathcal C^G}(\widetilde F,\widetilde F)
 =\operatorname{Ext}^2(F,F)^G$ in characteristic zero.
Functoriality of trace gives the commuting square
\[
\begin{array}{ccc}
 \operatorname{Ext}^2_{\mathcal C^G}(\widetilde F,\widetilde F)
 &\xrightarrow{\ \sigma_{\widetilde F}^2\ }&HH_{-2}(\mathcal C^G)\\
 \big\Vert&&\big\downarrow\scriptstyle{\mathrm{Forg}_*\;\simeq}\\
 \operatorname{Ext}^2(F,F)^G
 &\xrightarrow{\ \sigma_F^2\ }&HH_{-2}(A).
\end{array}
 \tag{F1R5.21}
\]
It identifies the two kernels exactly, as in HD-N12.
Equations (F1R5.15) and (F1R5.18) give a nonzero element
of that kernel, proving failure of weak $G$-semiregularity.

\paragraph{The one-summand extension.}
The conclusion already holds for an actual perfect complex
\[
 F=F'\oplus j_*E[b]
 \tag{F1R5.22}
\]
satisfying (F1R5.2), where $F'$ is arbitrary, $j:B+a\hookrightarrow A$
is any positive-dimensional proper abelian translate, and
$E$ is a coherent sheaf of positive generic rank on $B+a$.
Indeed paragraph1 only needs $H_1(B,\mathbf R)$ to be nonzero
and proper. Put $c=2n-\dim B$, so $1\le c\le2n-1$.
In (F1R5.8)--(F1R5.9) use $c$ normal covectors, and in
(F1R5.16)--(F1R5.17) the first nonzero character is
$\operatorname{ch}_c(j_*E)=r_B[B]$ and its detected component
is $r_BD_\kappa[B]\in H^{c+1}(A,\Omega_A^{c-1})$.
All remaining maps and the invariant kernel are unchanged.
This establishes (F1R5.22) directly, including arbitrary
additional summands; it supplies no replacement for the
actual direct-summand requirement.

\paragraph{Scope and precedence.}
This is an unconditional exclusion of the actual direct sums
(F1R5.1), with the full character (F1R5.2); it does not exclude
new cones or extensions with that character. The related
first-order direct-summand argument and structure-sheaf special
cases appear in Bhattacharjee--Bhattacharya,
\emph{Secant Objects for Weil Classes on Abelian Varieties},
\url{https://www.preprints.org/manuscript/202611.0672},
Lemma4 and Corollary3. Those unrefereed statements are comparison
sources only. Here the support movement, arbitrary coherent
sheaf, rational $B$-field, and every point-action kernel are
proved by the displayed maps. No originality claim is made.

\section{HD-F1-R06: the exact point-action test for linear averaging}

Let $f:A\to A$ be an origin-preserving automorphism and put
$\Psi=t_x^*f^*(-)\otimes P$ with $P\in\operatorname{Pic}^0(A)$.
On $HH^1(A)=H^0(A,T_A)\oplus H^1(A,\mathcal O_A)$ its
actual action is
\[
          \rho(\Psi)=\operatorname{diag}((df)^{-1},f^*).
 \tag{F1R6.1}
\]
Every point functor has identity action. Conversely,
$\rho(\Psi)=1$ forces $df=1$; lifting the group morphism
$f$ to its complex linear map on the universal cover gives
$f=1$. Thus a nontrivial linear automorphism cannot itself
be a point functor.

For any actual Fourier--Mukai equivalence $\Phi$, naturality
gives the exact conjugation relation
\[
 \rho(\Phi\Psi\Phi^{-1})
       =\Phi^{HH}\rho(\Psi)(\Phi^{HH})^{-1}.
 \tag{F1R6.2}
\]
Conjugation cannot change a nonidentity map into the identity.
Thus Fourier--Mukai conjugation cannot convert a finite
linear permutation of distinct graph tangent planes into
the point action required by S-14. On the support itself,
the point action sends $B+a$ to $B+a-x$ and preserves its
tangent plane, while tensoring changes no support.
For a commuting isogeny change $u$ and a lifted linear
automorphism $f'$, the exact equality
$df'=(du)^{-1}df\,du$ gives the same conclusion.

There is also an exact criterion for a fixed rational linear
map $f$ to act fibrewise on every member of a neighborhood
in the polarized Weil family. Such an action requires
$[f,I_s]=0$ for every period $I_s$ of paragraph1. Differentiating
and using $[f,I]=0$ gives $[f,P_Z]=0$ for every $Z$.
The matrix-unit calculation proves that $f$ commutes with
$\operatorname{Mat}_{2n}(\mathbf C)$, whose real commutant is
the scalar algebra $\mathbf C=\mathbf R[J]$. Its intersection
with $\operatorname{End}_{\mathbf Q}(H_1(A,\mathbf Q))$ is
exactly $K$: in a $K$-basis, a scalar preserving the rational
space must send the first basis vector into $K^{2n}$, so that
scalar belongs to $K$. Conversely every element of $K$
commutes with every allowed $I_s$. For a given integral
automorphism of the fixed lattice this proves
\[
 f\text{ acts fibrewise throughout the local Weil family}
                  \quad\Longleftrightarrow\quad f\in K^*.
 \tag{F1R6.3}
\]
Finite-order such maps are the roots of unity in $K$ which
preserve that lattice. This is a criterion for a fixed action
over the unchanged base, not for automorphisms which also
move the base point.

On the balanced product over
$K=\mathbf Q(i)$, the fixed integral $K$-action
$J=\operatorname{diag}(i1_n,-i1_n)$ is an automorphism
throughout the $K$-linear family. It sends the graph
$\{(z,Mz)\}$ to the graph $\{(z,-Mz)\}$, since
$z'=iz$ gives $-iMz=-Mz'$. Its action (F1R6.1) remains
nonidentity, so this actual family action still cannot
be substituted for S-14's point action through (F1R6.2).
The broader equivariant category for a linear group is
not assigned S-14's point-sector trace comparison here.
